\section{Результаты}

Робот запускался на корпусе \enquote{кино и сериалы} с двух источников. Прогон
ограничивался четырьмястами сохранёнными документами: цель прогона --- измерить
скорость и проверить работу всех механизмов, а не набрать полный корпус.
Итоги сведены в таблицу~\ref{tab:run}, разбивка по источникам --- в
таблицу~\ref{tab:sources}.

\subsection*{Что собрано}

\begin{table}[h]
\centering
\begin{tabular}{|l|r|}
\hline
\textbf{Показатель} & \textbf{Значение} \\
\hline
Сохранено документов              &       400 \\
Обойдено URL                      &       933 \\
Осталось в очереди                &   3\,221 \\
Объём \enquote{сырых} данных      & 81,2 МБ \\
Объём выделенного текста          &  3,02 МБ \\
Доля полезного текста             &  3,7\,\% \\
Средний размер документа          & 207,9 КБ \\
Средний объём текста              & 7\,918 символов \\
Размер файла базы                 & 87,6 МБ \\
\hline
\end{tabular}
\caption{Итоги прогона робота}
\label{tab:run}
\end{table}

Распределение объёма текста сильно скошено: медиана 4\,376 символов при среднем
7\,918, девяностый перцентиль 18\,934, максимум 81\,380. Это согласуется с
оценкой распределения, полученной в первой лабораторной работе.

\begin{table}[h]
\centering
\begin{tabular}{|l|r|r|r|l|}
\hline
\textbf{Источник} & \textbf{Док.} & \textbf{Сырой, КБ} & \textbf{Текст, симв.} & \textbf{Кодировка} \\
\hline
ru.wikipedia.org  & 383 & 212,2 &  7\,326 & UTF-8 \\
www.kino-teatr.ru &  17 & 109,8 & 21\,240 & windows-1251 \\
\hline
\end{tabular}
\caption{Документы по источникам}
\label{tab:sources}
\end{table}

Определение кодировки отработало: 383 документа сохранены как UTF-8, 17 ---
распознаны как windows-1251 и приведены к единому виду.

\subsection*{Скорость}

При паузе в одну секунду робот выдаёт около 0,95 страницы в секунду. Чередование
источников почти удваивает пропускную способность: пока идёт запрос к одному
хосту, пауза для второго истекает, поэтому суммарная скорость близка к одному
запросу в секунду, а не к одному запросу в две секунды.

Из этого следует оценка для полного корпуса. На 30 тысяч документов при таком
темпе уйдёт порядка 9--10 часов, и это без учёта страниц, которые скачиваются
ради ссылок, но документами не становятся, --- в нашем прогоне таких оказалось
больше половины (933 обойдённых URL против 400 сохранённых). Реальная оценка ---
около суток непрерывной работы. Робота можно останавливать и продолжать, так что
это не препятствие, но однопоточность --- главный ограничитель.

\subsection*{Блокировка второго источника}

Главный результат прогона отрицательный, и его стоит описать подробно.

Из 119 обойдённых адресов документов Кино-Театр.Ру сохранились лишь 17. Остальные
102 вернули не документ, а страницу одинакового размера в 1041 символ с текстом:

\begin{verbatim}
Кино-Театр.Ру Роботам запрещено посещать наш сайт.
Пожалуйста, подтвердите, что вы не робот. Проверка: введите...
\end{verbatim}

Сайт закрыт сервисом ddos-guard, и после примерно 470 запросов с паузой в одну
секунду защита перевела робота на капчу. Риск был отмечен ещё в первой
лабораторной работе, но тогда проверка пятью запросами прошла успешно и масштаба
проблемы не показала.

Неприятнее всего то, что робот не заметил подмены. Формально ответ приходил с
кодом 200, и страница с капчой --- такой же документ, как любой другой. От
попадания в базу четырёхсот одинаковых заглушек спас случай: порог
\texttt{min\_text\_length} в 1500 символов отсёк их как слишком короткие. Если бы
порог был ниже, корпус оказался бы испорчен незаметно.

По итогам в робота добавлены две вещи:
\begin{itemize}
    \item \texttt{blocked\_marker} --- признак страницы-заглушки в конфиге
          источника. Страница с этим признаком не сохраняется, а URL помечается
          в очереди как \texttt{failed} с указанием причины;
    \item своя пауза для источника: у Кино-Театр.Ру она увеличена до 5 секунд,
          что соответствует директиве \texttt{Crawl-delay} в его
          \texttt{robots.txt}.
\end{itemize}

Проверка показала, что распознавание работает: повторный запуск честно сообщает
\texttt{[блокировка]} вместо того, чтобы молча портить базу. Но сама блокировка
на момент написания отчёта не снята --- сайт отдаёт капчу даже на страницу
каталога. Вопрос остаётся открытым: нужно либо выждать и качать существенно
медленнее, либо заменить второй источник на более терпимый к роботам.

\pagebreak
